%! The Idea of Elimination — Unit 2, Lesson 2.1 — Group Activity
\documentclass[10pt]{article}
\usepackage{linalg-article}
\usepackage{linalg-boxes}
\newcommand{\TallMath}[1]{$\displaystyle #1\rule[-1.4em]{0pt}{3.2em}$}
\newcommand{\xx}{\mathbf{x}}
\newcommand{\bb}{\mathbf{b}}

\begin{document}

\pageheader{Unit 2, Lesson 2.1}{Group Activity}

\vspace{0.10in}

\begin{headlinebox}{blush}
\small\textbf{Elimination Assembly Line.} Every system today follows the same line: pick a
\textbf{pivot}, compute the \textbf{multiplier} $\ell=$ entry $\div$ pivot, subtract $\ell\times$
the pivot row to make a zero, reach \textbf{upper-triangular} form, and \textbf{back-substitute}.
Work with your partner and \emph{show every step}.
\end{headlinebox}

\vspace{0.10in}

\begin{tcolorbox}[colback=white, colframe=black!40, title=\textbf{Tier R --- Remediate},
    fonttitle=\bfseries, arc=1.5mm, left=3mm, right=3mm, top=2mm, bottom=2mm, breakable]
Solve $\ 2x + y = 7,\quad 4x + 3y = 17$. \; The pivot is $2$ and the multiplier is $\ell=4\div2=2$.
\begin{enumerate}[leftmargin=1.7em, itemsep=8pt]
  \item Elimination step: replace row~2 by (row~2)$-2\,($row~1$)$. What equation is left? \writeline
  \item That gives $y$. \; $y=\blank{1.4cm}$
\begin{work}
  (4x+3y)-2(2x+y) &= 17-14 \\
                y &= 3
\end{work}
  \item Back-substitute into row~1 to find $x$. \; $x=\blank{1.4cm}$
\begin{work}
  2x+3 &= 7 \\
     x &= 2
\end{work}
  \item \textbf{Check:} rebuild $\bb$ as $x\begin{bsmallmatrix} 2\\4 \end{bsmallmatrix}+y\begin{bsmallmatrix} 1\\3 \end{bsmallmatrix}$.
\begin{work}
  2\begin{bsmallmatrix} 2\\4 \end{bsmallmatrix} + 3\begin{bsmallmatrix} 1\\3 \end{bsmallmatrix}
      &= \begin{bsmallmatrix} 7\\17 \end{bsmallmatrix} \ \ \text{\checkmark}
\end{work}
\end{enumerate}
\end{tcolorbox}

\begin{tcolorbox}[colback=white, colframe=black!40, title=\textbf{Tier A --- Approaching Proficiency},
    fonttitle=\bfseries, arc=1.5mm, left=3mm, right=3mm, top=2mm, bottom=2mm, breakable]
\textbf{(1) Snack packs.} Each \textbf{Pack A} holds $1$ granola bar and $3$ juice boxes; each
\textbf{Pack B} holds $2$ granola bars and $1$ juice box. An order must use exactly $9$ granola
bars and $17$ juice boxes.
\begin{enumerate}[leftmargin=1.7em, itemsep=8pt]
  \item Write it as a system (bars row, juice row) with unknowns $a,b$ (packs of each). \writeline
  \item Solve by elimination. Choose the pivot so the multiplier is a whole number.
\begin{work}
  (3a+b)-3(a+2b) &= 17-27 \\
            -5b &= -10 \\
              b &= 2 \\
              a &= 9-4 = 5
\end{work}
  \item \emph{Interpret:} how many of each pack, and does the answer make sense? \writeline
\end{enumerate}

\vspace{4pt}
\tcbbreak
\textbf{(2) Three at once.} Solve by elimination (clear column~1, then column~2):
\[
  \begin{array}{r@{}c@{}l}
    x &{}+ y + z ={}& 6 \\
    2x &{}+ 3y + 4z ={}& 20 \\
    x &{}+ 4y + 9z ={}& 36
  \end{array}
\]
\begin{work}
  \ell_{21}=2:\quad y+2z &= 8 \\
  \ell_{31}=1:\quad 3y+8z &= 30 \\
  \ell_{32}=3:\quad 2z &= 6 \\
                     z &= 3 \\
                     y &= 2 \\
                     x &= 1
\end{work}
\end{tcolorbox}

\begin{tcolorbox}[colback=white, colframe=black!40, title=\textbf{Tier E --- Extension},
    fonttitle=\bfseries, arc=1.5mm, left=3mm, right=3mm, top=2mm, bottom=2mm, breakable]
\textbf{(1) Zero pivot.} In $\ 0x + 2y = 6,\ \ 4x + 3y = 17$, elimination can't start. Explain
why, fix it with a \textbf{row exchange}, then solve.
\begin{work}
  \text{swap rows:}\quad 4x+3y &= 17 \\
                          2y &= 6 \\
                           y &= 3 \\
                           x &= 2
\end{work}
\writelines{2}

\vspace{4pt}
\textbf{(2) Read the breakdown.} Run elimination on each system (multiplier $2$) and use the last
line to decide \emph{no solution} or \emph{infinitely many} --- and justify.
\begin{enumerate}[leftmargin=1.7em, itemsep=8pt]
  \item $\begin{aligned}[t] x + 2y &= 4 \\ 2x + 4y &= 9 \end{aligned}$
        \hspace{1.5em} No, or infinitely many? Why?
\begin{work}
  (2x+4y)-2(x+2y) &= 9-8 \\
                0 &= 1
\end{work}
        \writelines{2}
  \item $\begin{aligned}[t] x + 2y &= 4 \\ 2x + 4y &= 8 \end{aligned}$
        \hspace{1.5em} No, or infinitely many? Why?
\begin{work}
  (2x+4y)-2(x+2y) &= 8-8 \\
                0 &= 0
\end{work}
        \writelines{2}
\end{enumerate}
\end{tcolorbox}

\end{document}
